\section{Adapted OFDG: damping without changing the cell mean}
\label{chap:ofdg}
The DG state from \cref{chap:background} may contain unresolved oscillations.
Adapted OFDG reduces nonconstant content within each element while preserving
its mean. Three ingredients determine the operation: nested projections separate
orders of variation, interelement jumps measure roughness, and wave-speed-scaled
rates determine how strongly each part is damped.

\subsection{Relation to the published methods}
\label{sec:ofdg-relation-to-literature}
The name \emph{adapted OFDG} is used throughout this report for the specific
hybrid construction implemented here. Projection damping originates in
\cite{LuEtAl2021,LiuEtAl2022}; global scale normalization and exponential decay
are motivated by \cite{PengEtAl2024}. The earlier multidimensional OFDG
sensors evaluate vertex jumps, and the system formulation uses characteristic
variables. Here we average absolute jumps over faces and form separate rates
for physical conservative components. KXRCF gating is optional.

These choices define an implementation variant, not a literal reproduction of
one paper. The conservation and damping identities below follow from projection
algebra. Published convergence theorems for different sensors are not a proof
for this variant, particularly on curved elements with projected derivatives.

\subsection{Projections and a small example}
\label{sec:ofdg-mathematics}
For the nested mapped spaces from \cref{eq:ofdg-space}, let $\Pi_K^r$ be the
physical $L^2$ projection into $V^r(K)$:
\begin{equation}
 (\Pi_K^r w-w,v)_K=0\quad\hbox{for every }v\in V^r(K),
 \qquad (p,q)_K=\int_K pq\,dx.
 \label{eq:ofdg-l2-projection}
\end{equation}
This is the closest approximation in the integral squared-error norm.
Because constants belong to every space, the projection preserves the element
integral. In particular $\Pi_K^0w=|K|^{-1}\int_Kw\,dx$ is the constant cell mean.

For example, on $[-1,1]$ write
\begin{equation}
 w(x)=a_0+a_1x+a_2(x^2-1/3).
\end{equation}
The last two terms have zero mean. Reducing their coefficients changes the
slope and curvature but leaves $a_0$ unchanged. This orthogonal example is an
explanation, not a required storage basis: the code assembles projections in
MFEM's existing scalar finite element basis.

\Cref{fig:projection-damping} illustrates how attenuating zero-mean terms
changes a profile while preserving its average.
\begin{figure}[htbp]
\centering
\begin{tikzpicture}[x=4.3cm,y=1.65cm,>=Latex,font=\small]
 \draw[->,black!55] (-1.12,0)--(1.16,0) node[right] {$x$};
 \draw[->,black!55] (-1.08,0)--(-1.08,2.22) node[above] {$w$};
 \foreach \x in {-1,0,1} {
  \draw (\x,.025)--(\x,-.025) node[below] {$\x$};
 }
 \foreach \y in {0,1,2} {
  \draw (-1.095,\y)--(-1.065,\y);
  \node[left] at (-1.10,\y) {$\y$};
 }
 \draw[black!60,densely dotted,thick] (-1,1)--(1,1);
 \draw[blue!70!black,very thick,domain=-1:1,samples=81]
   plot (\x,{1+.55*\x+.7*(\x*\x-1/3)});
 \draw[orange!85!black,very thick,dashed,domain=-1:1,samples=81]
   plot (\x,{1+.33*\x+.175*(\x*\x-1/3)});
 \draw[blue!70!black,very thick] (-.90,2.18)--(-.73,2.18);
 \node[right] at (-.70,2.18) {before damping};
 \draw[orange!85!black,very thick,dashed] (.02,2.18)--(.19,2.18);
 \node[right] at (.22,2.18) {after damping};
 \node[fill=white,inner sep=2pt,anchor=south west] at (-.98,1.02)
       {shared mean $\overline w=1$};
\end{tikzpicture}
\caption{An analytic illustration of projection damping on $[-1,1]$.
The solid curve is $w=1+0.55x+0.7(x^2-1/3)$. Multiplying its linear and
quadratic shells by $0.6$ and $0.25$ gives the dashed curve. Both nonconstant
basis functions have zero integral, so the cell mean stays at one even though
the profile changes. The factors are chosen for illustration; these are not
simulation results or a positivity guarantee.}
\label{fig:projection-damping}
\end{figure}


Define the remainders and shells
\begin{align}
 \mathcal R_K^{(l)}&=I-\Pi_K^{l-1},\quad \Pi_K^{-1}:=\Pi_K^0,
 \label{eq:ofdg-high-pass}\\
 \mathcal Q_K^{(0)}&=\Pi_K^0,\qquad
 \mathcal Q_K^{(j)}=\Pi_K^j-\Pi_K^{j-1},\quad 1\leq j\leq k.
 \label{eq:ofdg-modal-shells}
\end{align}
Here $I$ is the identity operator. In particular,
$\mathcal R_K^{(0)}=\mathcal R_K^{(1)}=I-\Pi_K^0$: both remove the
constant mean and retain everything else. For $l\geq2$, the remainder
removes the whole order-$(l-1)$ subspace. The special convention for
$\Pi_K^{-1}$ ensures that even the solution-jump contribution cannot damp
the mean. A shell contains the variation added when increasing the reference-space order
from $j-1$ to $j$. These shells are mutually orthogonal in the physical inner
product. On tensor-product or curved elements, ``shell $j$'' means this nested
space difference; it does not mean total physical polynomial degree $j$.

\subsection{Projection matrices in the solution basis}
\label{sec:ofdg-projection-coefficients}
To apply these function-space operations to the stored DG coefficients, we
need matrices in the existing solution basis. Write
$u_h=\sum_{i=1}^{N_k}U_i\phi_i$ on $K$, so $U\in\mathbb R^{N_k}$
is a coefficient vector, not a list of point values in general.
The physical inner product in \cref{eq:ofdg-l2-projection} supplies the
matrix entries below; its evaluation on curved elements is derived in
\cref{eq:curved-weighted-masses}.

Let $\{\phi_i\}_{i=1}^{N_k}$ be the local degree-$k$ basis and
$\{\psi_a^{(r)}\}_{a=1}^{N_r}$ the degree-$r$ basis.  The code assembles
\begin{align}
  (M_k)_{ij}&=(\phi_i,\phi_j)_K,
  & (M_r)_{ab}&=(\psi_a^{(r)},\psi_b^{(r)})_K,\\
  (B_r)_{ai}&=(\psi_a^{(r)},\phi_i)_K.
\end{align}
Here $M_k\in\mathbb R^{N_k\times N_k}$ and
$M_r\in\mathbb R^{N_r\times N_r}$ are mass matrices, while
$B_r\in\mathbb R^{N_r\times N_k}$ transfers the inner products with
$u_h$ to the lower-order test basis. Testing the projection equation gives
$M_r C_r=B_r U$. The lower-order coefficients $C_r$ represent
$\Pi_K^r u_h$; lifting this same function to the original basis gives
$M_k U_r=B_r^{\mathsf T}C_r$. Eliminating $C_r$ yields the projector,
and subtraction from $U$ yields the remainder.
The matrix representation of $I-\Pi_K^r$ in the degree-$k$ coefficient basis is
\begin{equation}
  Q_r=M_k-B_r^\mathsf{T}M_r^{-1}B_r,
  \qquad
  S_r=M_k^{-1}Q_r
     =I-M_k^{-1}B_r^\mathsf{T}M_r^{-1}B_r.
  \label{eq:ofdg-projection-matrices}
\end{equation}
Thus applying $S_r$ to a degree-$k$ coefficient vector directly returns the
coefficients of the remainder $w-\Pi_K^r w$.  Here $I$ is the $N_k\times N_k$ identity matrix, and inverse notation denotes linear solves.  The code stores $S_r$
for $0\leq r\leq k$.

Crucially, this construction never converts the solution to Legendre or other
orthogonal modal coefficients.  It represents the same projection operator in
whatever supported scalar polynomial DG basis the host framework supplies.
Consequently, a nodal basis, a modal basis, or another polynomial coefficient
basis produces the same function-space operation up to quadrature and roundoff.
This basis-agnostic realization is what makes the implementation portable: a
second framework need not add a special OFDG solution basis, provided it can
form the corresponding polynomial-space projections.

On an affine element, all three mass
and coupling matrices contain the same constant physical Jacobian factor, which
cancels in $S_r$.  Therefore affine projection matrices are reused per finite-element signature.
For curved elements every mass and coupling matrix includes the physical
Jacobian weight at each quadrature point.  The lower-order spaces are mapped
reference polynomial spaces; their physical $L^2$ projectors remain nested.

For example, $S_0U$ is the nonconstant remainder, and
$(S_{j-1}-S_j)U$ represents shell $j$. Since $S_k=0$, these stored
remainders suffice for every shell in \cref{eq:ofdg-modal-shells}.

\subsection{Which derivatives enter the sensor?}
\label{sec:projected-derivatives}
First distinguish an exact derivative from the projected operation used in
its numerical approximation. A \emph{multi-index} is a tuple of nonnegative
integers. For spatial dimension $d$, define
\begin{equation}
 \alpha=(\alpha_1,\ldots,\alpha_d)\in\mathbb N_0^d,\qquad
 |\alpha|=\sum_{j=1}^d\alpha_j,\qquad
 \partial^\alpha=\partial_{x_1}^{\alpha_1}\cdots
                         \partial_{x_d}^{\alpha_d}.
 \label{eq:derivative-multi-index}
\end{equation}
Here $\mathbb N_0=\{0,1,2,\ldots\}$, $\alpha_j$ counts differentiations in
coordinate $x_j$, and $|\alpha|$ is their total order. For example, in two
dimensions $(1,0)$ means $\partial_x$, $(0,1)$ means $\partial_y$,
$(1,1)$ means a mixed second derivative, and $(2,1)$ means two
$x$-derivatives and one $y$-derivative. The zero multi-index means no
differentiation: $\partial^0v=v$. For sufficiently smooth functions within
an element, exact mixed partial derivatives commute, including on curved
elements. This statement concerns the physical function, not its discrete
approximation.

Define the projected first-derivative operator by
\begin{equation}
 \mathscr D_{K,j}v=\Pi_K^k(\partial_{x_j}v).
 \label{eq:projected-first-derivative}
\end{equation}
It first differentiates in physical coordinate $x_j$, then projects the result
back into $V^k(K)$. For repeated application, this report uses the explicit
convention
\begin{equation}
 \mathscr D_K^\alpha
 =\mathscr D_{K,d}^{\alpha_d}\cdots
   \mathscr D_{K,2}^{\alpha_2}\mathscr D_{K,1}^{\alpha_1},
 \qquad \mathscr D_K^0=I.
 \label{eq:projected-multi-index-order}
\end{equation}
Operators act from right to left: all $x_1$ operations come first, then
$x_2$, and so on. This is the sequence selected by the implementation.
Thus $\mathscr D_K^{(1,1)}v=\mathscr D_{K,y}\mathscr D_{K,x}v$ in 2D;
the superscript labels that particular sequence, not all its permutations.

On affine elements whose reference family is closed under differentiation,
these operators recover the exact polynomial derivatives up to numerical
precision, and their order is immaterial. On a curved element they represent
\emph{successively projected} physical derivatives. Inserting a projection
between differentiations can destroy commutation. Consequently
$\mathscr D_{K,y}\mathscr D_{K,x}v$ and
$\mathscr D_{K,x}\mathscr D_{K,y}v$ need not agree, even though
$\partial_y\partial_xv=\partial_x\partial_yv$.
\Cref{sec:curved-recursive-derivatives} derives the difference and explains
what the multi-index storage does and does not imply.

\subsection{From reference derivatives to physical coefficients}
\label{sec:ofdg-coefficient-derivatives}

The sensor needs derivative traces on both sides of a face. We first
construct derivative coefficients within each element, then evaluate their
traces. This avoids requesting higher basis derivatives from MFEM.
Let $\widehat\phi_i=\phi_i\circ T_K$ be the reference basis and let
$\{\boldsymbol\xi_q\}_{q=1}^{N_k}$ be the finite element's interpolation
points. These are distinct from the integration points used for volume or
face quadrature. The supported element must provide an invertible evaluation
matrix at these points. For each reference direction $r\in\{1,\ldots,d\}$,
define the two $N_k\times N_k$ matrices
\begin{equation}
  V_{qi}=\widehat\phi_i(\boldsymbol{\xi}_q),
  \qquad
  G_{qi}^{(r)}
  =\frac{\partial\widehat\phi_i}{\partial\xi_r}(\boldsymbol{\xi}_q).
\end{equation}
Rows index sample points and columns index basis functions. Thus
$V\boldsymbol u$ gives solution values at the interpolation points, while
$G^{(r)}\boldsymbol u$ contains
reference-derivative values at the interpolation points.  To express those derivative values as coefficients, solve
$V\boldsymbol u_{\xi_r}=G^{(r)}\boldsymbol u$. Because reference
polynomial differentiation stays in the supported space, this reconstructs
the derivative exactly up to roundoff. Therefore
\begin{equation}
  \boldsymbol{u}_{\xi_r}=D_{\xi_r}\boldsymbol{u},
  \qquad
  D_{\xi_r}=V^{-1}G^{(r)}.
  \label{eq:ofdg-reference-derivative-matrix}
\end{equation}
The code's \texttt{AssembleReferenceDerivatives} builds $V$ with MFEM's
\texttt{CalcShape}, builds $G^{(r)}$ with \texttt{CalcDShape}, and
explicitly forms $V^{-1}G^{(r)}$. For a Lagrange nodal basis at its own nodes,
$V=I$. Keeping $V$ explicit also handles supported bases whose coefficients
are not those nodal values; invertibility is essential.

For a small example, use $(1,\xi,\xi^2)$ on $[0,1]$ and sample at
$0,\tfrac12,1$. Then
\begin{equation}
 V=\begin{pmatrix}1&0&0\\1&1/2&1/4\\1&1&1\end{pmatrix},\quad
 G=\begin{pmatrix}0&1&0\\0&1&1\\0&1&2\end{pmatrix},\quad
 D_\xi=\begin{pmatrix}0&1&0\\0&0&2\\0&0&0\end{pmatrix}.
\end{equation}
For $\boldsymbol u=(a,b,c)^{\mathsf T}$, $G\boldsymbol u$ contains
the sampled slopes $(b,b+c,b+2c)^{\mathsf T}$, whereas
$D_\xi\boldsymbol u=(b,2c,0)^{\mathsf T}$ contains the coefficients
of $b+2c\xi$. This is why values and coefficients cannot generally be
interchanged.

Now introduce the affine element map
\begin{equation}
 \boldsymbol x=T_K(\boldsymbol\xi)=\boldsymbol b_K+J_K\boldsymbol\xi,
 \qquad (J_K)_{jr}=\frac{\partial x_j}{\partial\xi_r}.
 \label{eq:ofdg-affine-map}
\end{equation}
The vector $\boldsymbol b_K$ translates the element and the invertible
$d\times d$ matrix $J_K$ describes scaling, rotation, and shear.
It is constant within this element but may differ between elements.
The chain rule $\partial_{x_j}=\sum_r(J_K^{-1})_{rj}\partial_{\xi_r}$
gives the physical coefficient matrices
\begin{equation}
  D_{x_j}^{K}
  =\sum_{r=1}^d(J_K^{-1})_{rj}D_{\xi_r}.
  \label{eq:ofdg-physical-derivative-matrix}
\end{equation}
The inverse Jacobian $J_K^{-1}$ is read once at the reference-element centre and
one set of matrices $\{D_{x_j}^{K}\}_{j=1}^d$ is stored for every element.

For a curved map $J_K$ varies within the element, so the constant
combination in \cref{eq:ofdg-physical-derivative-matrix} is insufficient.
The matrix representing \cref{eq:projected-first-derivative} is instead
$D_{K,j}=(M_k^K)^{-1}G_{K,j}$, with the physical integral matrix
$G_{K,j}$ derived in \cref{eq:curved-derivative-assembly}.
Unlike $G^{(r)}$, which stores pointwise reference derivative values,
$G_{K,j}$ stores physical inner products. Repeated application follows
\cref{eq:projected-multi-index-order}; the parent rule and the projection
error are developed in \cref{sec:curved-recursive-derivatives}.

\paragraph{Coordinate derivatives and normal derivatives.}
A physical coordinate derivative differentiates along $x_j$. At a face
point $x_q$ with outward unit normal
$\boldsymbol n_K(x_q)=(n_1(x_q),\ldots,n_d(x_q))$, the normal derivative is
\begin{equation}
 \partial_{n_K}u_h(x_q)=\sum_{j=1}^d n_j(x_q)\partial_{x_j}u_h(x_q).
 \label{eq:ofdg-normal-derivative}
\end{equation}
On an affine planar face it can be evaluated from
$\sum_j n_jD_{x_j}^K\boldsymbol u$ and the face basis values. On a curved
face, contracting the projected coordinate derivatives gives an approximation
$\sum_j n_j(x_q)(\mathscr D_{K,j}u_h)(x_q)$; the normal must be evaluated
at that point. This need not equal the exact normal derivative or its single
physical projection. The implemented OFDG sensor uses jumps of the
\emph{coordinate derivatives} in \cref{eq:ofdg-face-jump}, not jumps of
this contraction. Normals instead enter its wave speeds and KXRCF's inflow
selection. No separate normal-derivative matrix is required by these sensors.

Only first reference derivatives of the supplied finite-element basis are
requested from MFEM.  Repeated application of the resulting coefficient
operators generates every higher derivative needed by OFDG.  In framework
terms, the construction requires scalar polynomial DG basis values and first
derivatives, quadrature, nested lower-order polynomial spaces or equivalent
projection operators, element and face transformations, and local
degree-of-freedom access.  It therefore avoids a modal-basis requirement, but it
does not claim support for arbitrary nonpolynomial, $H(\mathrm{div})$,
$H(\mathrm{curl})$, or mixed finite-element spaces.

All multi-indices with total degree at most $k$ are organized in a directed
acyclic graph.  Each nonzero multi-index selects one parent of degree one less and
one physical derivative direction.  Walking this graph evaluates every required
derivative once.  The number of stored derivative states per element and component
is
\begin{equation}
  \sum_{l=0}^k\binom{l+d-1}{d-1}=\binom{k+d}{d}.
\end{equation}

\subsection{Derivative jumps and their face evaluation}
\label{sec:ofdg-face-evaluation}
For an interior face $F=\partial K^-\cap\partial K^+$, let $w^-$ and
$w^+$ be traces from its two adjacent elements. The jump is their difference.
For component $c\in\{1,\ldots,m\}$ and derivative order
$l\in\{0,\ldots,k\}$, the sensor adds the mean absolute jumps of all
projected derivatives of total order $l$. With $|F|$ the physical face
measure, these definitions read
\begin{equation}
 [\![w]\!]_F=w^-|_F-w^+|_F,
 \qquad
 J_{F,c}^{(l)}=\sum_{|\alpha|=l}\frac1{|F|}\int_F
 \left|[\![\mathscr D^\alpha u_{h,c}]\!]_F\right|\,ds.
 \label{eq:ofdg-face-jump}
\end{equation}
The sum runs over
all nonnegative multi-indices with $|\alpha|=l$, once each. In 2D, the
$l=1$ sum includes $(1,0)$ and $(0,1)$; the $l=2$ sum includes $(2,0)$,
$(1,1)$, and $(0,2)$. The mixed term is counted once, using
\cref{eq:projected-multi-index-order}; there is no additional reversed-order
term or multiplicity factor. For $l=0$ the only index is $(0,0)$, so the
sensor measures jumps of the solution itself. Each trace uses its own
element's derivative operator at the same physical face point. Face measure $|F|$ is
length in 2D or area in 3D; a 1D point face uses the ordinary absolute jump.
Taking the absolute value before integration prevents cancellation. For
constant traces 2 and 1.25, $J_F^{(0)}=0.75$ regardless of face size.


MFEM remains the sole source of mesh connectivity.  The setup loops over MFEM
face numbers and obtains the two adjacent elements from MFEM's face-transformation
data.  Boundary faces are skipped.  For each
interior affine face, a quadrature rule of order $2k+1$ is selected.  Curved
faces use an overintegrated rule that also depends on the geometry order.  Let $x_q$ be a physical face quadrature point and let
$\boldsymbol\xi_{F,q}^{\pm}=T_{K^{\pm}}^{-1}(x_q)$ be its reference
coordinates in the two adjacent elements. If their coefficient counts are
$N_k^-$ and $N_k^+$, then $E_F^\pm$ has size $N_{q,F}\times N_k^\pm$.
Each row converts one side's coefficients to a trace value at $x_q$.
The following data are cached directly at the MFEM face index:
\begin{equation}
  (E_F^-)_{qi}=\widehat\phi_i(\boldsymbol{\xi}_{F,q}^-),
  \qquad
  (E_F^+)_{qi}=\widehat\phi_i(\boldsymbol{\xi}_{F,q}^+),
  \qquad
  \widehat{\omega}_q=\frac{\omega_q J_F(\xi_q)}{\sum_p\omega_p J_F(\xi_p)}.
  \label{eq:ofdg-face-evaluation-matrices}
\end{equation}
Here $J_F$ is the physical face Jacobian.  For affine faces it is constant and cancels between the
integral and $|F|$ in~\eqref{eq:ofdg-face-jump}; hence the normalized reference
weights $\widehat{\omega}_q$ are sufficient.  All derivative traces and all
components on one side of a face are evaluated together in one dense matrix
multiplication.  There is no face-pattern matching or deduplication.

If $\boldsymbol{u}_{\boldsymbol{\alpha}}^-$ and
$\boldsymbol{u}_{\boldsymbol{\alpha}}^+$ are the coefficient vectors of the
same physical derivative on the two neighboring elements, the actual quadrature
calculation is
\begin{equation}
  J_{F,c}^{(l)}
  \approx\sum_{|\boldsymbol{\alpha}|=l}\sum_{q=1}^{N_{q,F}}
  \widehat{\omega}_q
  \left|
    \bigl(E_F^-\boldsymbol{u}_{\boldsymbol{\alpha}}^-\bigr)_q
    -\bigl(E_F^+\boldsymbol{u}_{\boldsymbol{\alpha}}^+\bigr)_q
  \right|.
  \label{eq:ofdg-discrete-face-jump}
\end{equation}
The implementation therefore performs one subtraction and one absolute value for
each derivative, component, and quadrature point.  No squared-jump accumulator or
final square root is required.

Thus the full coefficient-to-sensor route is: differentiate the element
coefficients, multiply by $E_F^-$ and $E_F^+$, subtract matching traces,
take absolute values, and average with normalized physical face weights.
For curved elements, the pointwise normal and surface factor are derived in
\cref{eq:curved-face-metric,eq:curved-face-normalization}.

\begin{figure}[htbp]
\centering
\begin{tikzpicture}[node distance=6mm,>=Latex,
 box/.style={draw,rounded corners,fill=blue!4,text width=43mm,align=center,minimum height=12mm}]
 \node[box](traces){solution and projected derivative traces};
 \node[box,right=of traces](jumps){face-averaged absolute jumps};
 \node[box,below=of jumps](rates){rates from jumps, amplitude, size, and speed};
 \node[box,left=of rates](decay){local shell attenuation; mean unchanged};
 \draw[->](traces)--(jumps);\draw[->](jumps)--(rates);\draw[->](rates)--(decay);
\end{tikzpicture}
\caption{Adapted OFDG separates face-based sensing from element-local damping.
KXRCF, when enabled, selects which elements receive the final operation.}
\label{fig:ofdg-pipeline}
\end{figure}

\subsection{Damping rates and physical scales}
\label{sec:ofdg-rates}
Let $\overline u_c=|\Omega|^{-1}\int_\Omega u_{h,c}\,dx$ and define the ideal
component amplitude
\begin{equation}
 A_c=\|u_{h,c}-\overline u_c\|_{L^\infty(\Omega)}.
 \label{eq:ofdg-global-scale}
\end{equation}
The code approximates this amplitude from volume-quadrature samples and an
integrated mean; it does not optimize the continuous polynomial maximum.
Adapted OFDG sets its component rates to zero when the sampled amplitude is
at most $10^{-14}$. This absolute cutoff limits exact rescaling invariance
near vanishing amplitudes.

The one-sided speed $\beta_{K,F}$ is the largest sampled normal radius on
face $F$ using the state on side $K$. Prescribed advection uses
$|\boldsymbol a\cdot\boldsymbol n_K|$; Euler uses
$|\boldsymbol v_K\cdot\boldsymbol n_K|+c_K$ evaluated from face traces.
The other element can have a different radius. Normals are evaluated at the
same physical face quadrature points. Unit speed is a compatibility policy,
rather than the state-dependent physical-speed policy defined here.

With $k\geq1$ the solution-space order, $|K|$ and $|\widehat K|$ the
physical and reference element volumes, and
$h_K=(|K|/|\widehat K|)^{1/d}$, the rates are
\begin{equation}
 \delta_{K,c}^{(l)}=
 \frac{2l+1}{2(2k-1)l!}\frac{h_K^{l-1}}{A_c}
 \sum_{\substack{F\subset\partial K\\F\in\mathcal F_h^{\rm i}}}
 \beta_{K,F}J_{F,c}^{(l)},\qquad 0\leq l\leq k,
 \label{eq:ofdg-damping-rate}
\end{equation}
Here $\mathcal F_h^{\rm i}$ denotes the set of interior mesh faces,
$\partial K$ is the element boundary, and $l!$ is the factorial of the total
derivative order, with $0!=1$. The sum collects contributions from interior
faces of $K$ only. Physical boundary conditions
belong to the DG flux operator. Face and volume integrals in the implementation
use quadrature, including geometry-dependent weights on curved cells.

The rate has inverse-time units: the derivative jump contributes $u/$length$^l$,
$h_K^{l-1}$ contributes length$^{l-1}$, speed contributes length/time, and
$A_c$ divides out the state scale. At fixed geometry and amplitude, larger speeds or jumps strengthen
attenuation; the factor $h_K^{l-1}$ sets the length scaling for each derivative order. For $k=2$, $l=0$, $h_K=0.1$, $A=1$,
$\beta=2$, and one jump $J=0.3$, the contribution is
$(1/6)(0.1^{-1})(2)(0.3)=1$ per unit time.

With fixed prescribed speeds and away from the cutoff, the jump/amplitude
ratio is unchanged under $u_c\mapsto\lambda u_c+b$, $\lambda\ne0$.
Arbitrary component scaling or shifting of Euler or Burgers states can change
the physical speeds, so the complete nonlinear filtering map does not inherit
that statement unconditionally.

\subsection{Conservation, dissipation, and exact frozen decay}
\label{sec:ofdg-decay}
We now combine the two element-local building blocks. For each derivative
order $l$, the rate $\delta_{K,c}^{(l)}$ measures how strongly to act, and
the remainder $\mathcal R_K^{(l)}u_{h,c}$ selects which content may be
removed. Adding these contributions defines the damping operator
\begin{equation}
 \mathcal D_{K,c}u_{h,c}=\sum_{l=0}^k\delta_{K,c}^{(l)}
          \mathcal R_K^{(l)}u_{h,c}.
 \label{eq:ofdg-damping-operator}
\end{equation}
Every remainder has zero integral. For ideal physical projections,
\begin{align}
 \int_K\mathcal D_{K,c}u_{h,c}\,dx&=0,
 \label{eq:ofdg-local-conservation}\\
 (\mathcal D_{K,c}u_{h,c},u_{h,c})_K
 &=\sum_{l=0}^k\delta_{K,c}^{(l)}
       \|\mathcal R_K^{(l)}u_{h,c}\|_{L^2(K)}^2\geq0.
 \label{eq:ofdg-dissipation}
\end{align}
Thus subtracting $\mathcal D u$ removes nonconstant squared-norm content
without changing the cell mean. These are damping-step identities, not a
positivity or entropy-stability proof for the full nonlinear time update.
The numerical realization preserves its quadrature-weighted physical cell
means to roundoff; agreement with exact integrals also depends on quadrature.

Freeze the rates at the candidate state and solve
$\partial_s w+\mathcal D w=0$ for a pseudo-time interval $\tau$. Nested
projections give the explicit solution
\begin{equation}
 w(\tau)=\Pi_K^0u_{h,c}+\sum_{j=1}^k
  \exp\!\left(-\tau\sum_{l=0}^j\delta_{K,c}^{(l)}\right)
  (\Pi_K^j-\Pi_K^{j-1})u_{h,c}.
 \label{eq:ofdg-exact-decay}
\end{equation}
The mean is copied and each higher shell is multiplied by a factor between
zero and one. For order three, rates $l=0$ and $l=1$ affect all nonconstant
shells; rate $l=2$ affects shells two and three; rate $l=3$ affects shell three.
``Exact'' refers to this frozen damping equation. The rates are recomputed
from the evolving solution, and operator splitting remains part of the method.

The public stabilization routine returns coefficients of $\mathcal D u$ for
a caller that wants a semidiscrete term. The default post-step implementation uses the exponential
\texttt{CompDecay} map after a full RK4 step, with $\tau$ equal to the step
size. Basis-independent matrix construction is given in
\cref{sec:ofdg-projection-coefficients}.
